% Pursuing Stacks (cotes 134-2 to 134-8): archive cote 134-6, site page 78 (PDF page 79), typescript page 396.
% typescript. Draft compared: ps5.tex:5403-5493.
% AI-transcribed by Claude Opus 5.5 (claude-opus-5-5, adaptive thinking, effort high), Message Batches API, batch ps3 (msgbatch_01XdFmuVVBAcgc7EFxdtCqsj), 2026-11-12; not yet checked by eye.
\scanpage{134-6/78}
\mspage{396}

the category of \hand{$\mathcal{M}$}-valued \uline{sheaves} on the topos $B^{\wedge}$.

This equivalence, defined up to unique isomorphism, is deduced
from the diagram of canonical equivalences
\diagram{\begin{tikzcd}[column sep=small]
\underline{\mathrm{Hom}}_{\hand{!}}(A^{\wedge},\hand{$\mathcal{M}$})\hand{$\overset{\mathrm{def}}{=}(A^{\wedge})^{\mathcal{M}}$} \arrow[dr, "\hand{$\approx$}"] & & \hand{$(B^{\wedge})_{\mathcal{M}}\overset{\mathrm{def}}{=}$}\underline{\mathrm{Hom}}^{\hand{!}}(B^{\wedge\circ}\unsure{;|,}\hand{$\mathcal{M}$}) \arrow[dl, "\hand{$\approx$}"'] \\
(40) & \underline{\mathrm{Hom}}(A,\hand{$\mathcal{M}$}) = \underline{\mathrm{Hom}}(B^{\circ},\hand{$\mathcal{M}$}) & ,
\end{tikzcd}}

and depends upon the choice of a quasi-inverse of the second vertical
equivalence in (40) (which choice can be made, via the dual of (35),
via the choice of small inverse limits in \hand{$\mathcal{M}$}).

\uline{\unsure{Remark|Rgmark}.} It is felt that the way we got the equivalence (39) via (40),
the rôle of A and B in it is symmetric. To give a more precise statement,
consider the equivalence deduced from (39) by passing to the dual catego-%
ries of the two members - using the tautological isomorphisms (6) of
page 385, we get an equivalence \deleted{\hand{(\illegible{1}}}
\[ (39')\quad \hand{$(\delta^{\mathcal{M}}_{A})$}^{\unsure{\circ|'}}:\ \underline{\mathrm{Hom}}^{\hand{!}}(A^{\wedge\circ},\hand{$\mathcal{N}$}) \hand{$\xrightarrow{\approx}$} \underline{\mathrm{Hom}}_{\hand{!}}(B^{\wedge},\hand{$\mathcal{N}$}) \quad , \quad \text{Where } \hand{$\mathcal{N}=\mathcal{M}^{\circ}$}\ ; \]
th\unsure{\over{i}{e}}s equivalence is \uline{canonically} \uline{quasi-inverse} \uline{to the equivalence} \xout{\illegible{1}}
\hand{$\delta^{\mathcal{N}}_{B}$} \uline{in opposite direction}, associated to the pair (B,\hand{$\mathcal{N}=\mathcal{M}^{\circ}$}) instead of
(A,\hand{$\mathcal{M}$}) .

In the rest of this subsection D), we'll elaborate on some particu-%
lar cases of the equ\unsure{\xout{m}}ivalence (39) between cosheaves and sheaves.

\inserted{Case}\deleted{(}1\textsuperscript{o}). Assume \hand{$\mathcal{M}$} = (Sets), then by prop.1 (p.384) the right hand
side of (39) is canonically equivalent to $B^{\wedge}$ itself, hence we get an
equivalence
\[ (41)\qquad B^{\wedge} \hand{$\xrightarrow{\approx}$} \underline{\mathrm{Hom}}_{\hand{!}}(A^{\wedge},\ (\mathrm{Sets})) \quad . \]
If we want to keep track of the symmetry aspect described in the remark
above, we may consider \xout{(41)xas} the functor (41) as being deduced from a
canonical"pairing" between the categories $A^{\wedge}$ and $B^{\wedge}$
\[ (42)\qquad \hand{$\delta_{A}$}:\ A^{\wedge}\ \mathrm{x}\ B^{\wedge} \hand{$\longrightarrow$} (\mathrm{Sets}) \hand{,} \]
which is an object in
\[ \underline{\mathrm{Hom}}_{\hand{!!}}(A^{\wedge},\ B^{\wedge};\ (\mathrm{Sets})) \]
i.e. which commutes to small direct limits in e ach variable. (For this
interpretation, compare the dual statement contained in \unsure{formula|fbrmula} (7) of
page 385 - and note that (41), being an equivalence, commutes to
small direct limits i.e. is in a category $\underline{\mathrm{Hom}}_{\hand{!}}(B^{\wedge},\ \underline{\mathrm{Hom}}_{\hand{!}}(\ \ ,\ \ ))$ .)

\margin{78}
